\begin{proposition}[Proved under stated hypotheses: a single positive-genus exceptional curve]
\label{D:single-curve-residue-tameness}
Let $k$ be algebraically closed, let $(S,s)$ be a normal surface
singularity of finite type over $k$, and suppose that a projective
resolution $Y\to S$ has reduced exceptional locus a single smooth
connected curve $C$, of genus $g>0$ and self-intersection $C^2=-d<0$.
Let $S'\to S$ be a finite generically separable morphism of normal
integral surfaces, \'etale away from the points over $s$, of degree $n$.
For the extensions $F_i$ of the divisorial valuation $C$, write
$e_i$ for the ramification index, $f_i$ for the residue-field degree,
$f_i^{\mathrm{sep}}$ for its separable degree, and $\delta_i$ for
the coefficient of the relative canonical divisor on smooth models.
Then
\[
 0\leq\sum_i f_i(\delta_i-e_i+1)
 \leq \frac{2g-2}{d}
             \left(n-\sum_i e_i f_i^{\mathrm{sep}}\right).
 \tag{D.PG.1}
\]
In particular the cover is tame at every valuation above $C$ if
every residue extension is separable. If $g=1$, all the residue
extensions are automatically separable and the cover is tame there.
No $\mathbb Q$-Gorenstein assumption on $S$ or $S'$ is required.
\end{proposition}

\begin{proof}
We shrink $S$ to an affine neighborhood of $s$ with no other
singularities. Normalize $Y$ in the function field of $S'$ and
resolve that normal surface by a morphism that is an isomorphism
off the inverse image of $C$, where the normalization is already
smooth. This gives a smooth projective
$S'$-surface $Z$ and a proper generically finite morphism
$f:Z\to Y$ of degree $n$. Resolution by normalized blowups is
\cite[Theorem 54.14.5, Tag 0BGP]{StacksDSurfaceResolution}.
We may make $Z\to S'$ minimal while keeping a morphism to $Y$.
Indeed, any exceptional smooth rational curve of self-intersection
$-1$ maps into $C$, and its image is a point: a nonconstant map
from $\mathbb P^1$ to a smooth curve of positive genus is impossible,
including for an inseparable map, by L\"uroth's theorem.
Contract this curve using
\cite[Lemma 54.16.8, Tag 0C2M]{StacksDSurfaceResolution}.

We spell out descent of $f$ through the contraction $b:Z\to Z_1$.
The scheme-theoretic image $W$ of $(b,f):Z\to Z_1\times_S Y$ is
integral and proper over $Z_1$. Away from the contracted point its
projection is an isomorphism, and over that point its underlying
fibre is a single point because $f$ is constant on the contracted
curve. Thus $W\to Z_1$ is proper and quasi-finite, hence finite.
It is birational, and $Z_1$ is normal, so it is an isomorphism.
This gives the required descended morphism. Each contraction lowers
the number of exceptional components, so the procedure terminates.

On this minimal $Z$, the canonical divisor has nonnegative
intersection with every exceptional integral curve $D$ over $S'$.
The exceptional intersection matrix is negative definite
\cite[Section 10.1, Theorem 10.1]{Kollar2013D}, hence $D^2<0$.
Adjunction gives
\[
                  K_Z\cdot D=2p_a(D)-2-D^2\geq0.
\]
If $p_a(D)=0$, then $D$ is a smooth rational curve; minimality
excludes $D^2=-1$, so the inequality still holds. We use here
the adjunction formula $(K_Z+D)\cdot D=2p_a(D)-2$ for a Cartier
curve on a smooth surface; it follows by applying degree to
$\omega_D=(\omega_Z\otimes\mathcal O_Z(D))|_D$.

Let the $F_i$ now denote the actual curves on $Z$ dominating $C$,
and let $Q_j$ be all the remaining exceptional curves. Their
images under $f$ are points of $C$. We have
\[
             f^*C=\sum_i e_iF_i+\sum_j m_jQ_j,
             \qquad m_j>0,\qquad \sum_i e_if_i=n.
 \tag{D.PG.2}
\]
The last equality is the degree formula for the finite normalization
over the discrete valuation ring at the generic point of $C$;
it includes the inseparable parts of the residue degrees.
By the projection formula and nonnegativity of intersections
between distinct integral curves on a smooth surface,
\[
 -df_i=(f^*C)\cdot F_i
       =e_iF_i^2+\sum_{h\ne i}e_h(F_h\cdot F_i)
                      +\sum_jm_j(Q_j\cdot F_i),
\]
so $e_i(-F_i^2)\geq df_i$.

Let $\widetilde F_i$ be the normalization of $F_i$.
Factoring $\widetilde F_i\to C$ into its purely inseparable
and separable parts, the purely inseparable part preserves the
genus over the perfect field $k$, by
\cite[Proposition 53.13.7 and Lemma 53.13.5, Tags 0CD2 and 0CD0]{StacksDSurfaceResolution}.
Riemann--Hurwitz for the
separable part gives
\[
 2p_a(F_i)-2\geq 2g(\widetilde F_i)-2
                   \geq f_i^{\mathrm{sep}}(2g-2).
\]
For the Riemann--Hurwitz statement we use
\cite[Lemma 53.12.2, Tag 0C1D]{StacksDSurfaceResolution}.
Together with nefness on the $Q_j$, adjunction, and (D.PG.2),
this yields the lower bound
\[
 \begin{split}
 K_Z\cdot f^*C
 &\geq\sum_i e_i\bigl(2p_a(F_i)-2-F_i^2\bigr)\\
 &\geq(2g-2)\sum_i e_i f_i^{\mathrm{sep}}+d\sum_i f_i.
 \end{split}
 \tag{D.PG.3}
\]

The determinant of the differential of $f$ defines an effective
relative canonical divisor $R_f$, with
$K_Z=f^*K_Y+R_f$. The morphism is \'etale outside $f^{-1}C$,
so $R_f$ is supported on the $F_i$ and $Q_j$, with coefficient
$\delta_i$ at $F_i$. As $(f^*C)\cdot Q_j=0$, the projection
formula and adjunction on $C$ give the exact identity
\[
 \begin{split}
 K_Z\cdot f^*C
  &= nK_Y\cdot C-d\sum_i f_i\delta_i\\
  &= n(2g-2+d)-d\sum_i f_i\delta_i.
 \end{split}
 \tag{D.PG.4}
\]
Combining (D.PG.3) and (D.PG.4), and using
$n=\sum_i e_i f_i$, proves the upper inequality in (D.PG.1).

For clarity, the lower inequality and its equality case include
inseparable residue extensions. Choose uniformizers $t,s$ at
$C,F_i$ and separating residue parameters $z,w$ over $k$.
At the corresponding discrete valuation ring write $t=u s^{e_i}$.
Then
\[
 d(us^{e_i})\wedge dz
   =s^{e_i-1}\bigl(e_i u\,ds\wedge dz+s\,du\wedge dz\bigr).
\]
The order of this differential relative to $ds\wedge dw$ is
$\delta_i$, so $\delta_i\geq e_i-1$. Equality holds exactly
when $e_i$ is nonzero in $k$ and the residue extension is
separable, which is precisely tame ramification. The determinant
is nonzero because the extension of the surface function fields
is separable. Thus every summand in (D.PG.1) is nonnegative.
If all $f_i=f_i^{\mathrm{sep}}$, its right-hand side is zero.
If $g=1$, it is zero without any separability assumption.
In either case all summands vanish, giving the conclusion.
\end{proof}

\begin{remark}[Scope of the higher-genus reduction]
For the single-curve singularities obtained by contracting a strict
vertical fibre in the mutation construction, the hypotheses hold
with $d=1$. Consequently a finite cover with wild ramification at
such an exceptional valuation must have an inseparable residue
extension at at least one valuation above it. For a Galois cover,
all residue separability types are conjugate, so every such wild
case has inseparable residue extension. We do not infer finite
generation, a finite image of wild inertia, or a uniform bound on
the orders of finite covers from (D.PG.1). In genus at least two
the remaining inseparable-residue case is still an explicit gate.
\end{remark}
